\specialsection{Multi-indices}

Let I be a set. For \(\alpha = (\alpha_{i})\) and \(\beta = (\beta_{i})\) belonging to \(\mathbf{N}^{(I)}\) , we set (cf. Alg., ch. III, \(3^{\mathrm{e}}\) ed., beginning of the chapter):

\[
| \alpha | = \sum_ {i \in I} \alpha_ {i}
\]

\[
\alpha ! = \prod_ {i \in I} \alpha_ {i}!
\]

\[
\alpha + \beta = (\alpha_ {i} + \beta_ {i}) _ {i \in I}
\]

\[
((\alpha , \beta)) = \prod_ {i \in I} \frac {(\alpha_ {i} + \beta_ {i}) !}{\alpha_ {i} ! \beta_ {i} !}.
\]

We set \(\alpha \leqslant \beta\) when \(\alpha_{i} \leqslant \beta_{i}\) for all \(i \in I\) , which is equivalent to the existence of \(\gamma \in \mathbf{N}^{(I)}\) such that \(\beta = \alpha + \gamma\) ; this element \(\gamma\) is then unique, and is denoted \(\beta - \alpha\) .

If \(\mathbf{x} = (x_i)_{i \in I}\) is a family of elements commuting pairwise in a ring A possessing an identity element 1, we set:

\[
\mathbf {x} ^ {\alpha} = \prod_ {i \in \mathbf {I}} x _ {i} ^ {\alpha_ {i}} \quad (\text { with   the   convention   that } x ^ {0} = 1).
\]

One has \(\mathbf{x}^{\alpha +\beta} = \mathbf{x}^{\alpha}\cdot \mathbf{x}^{\beta}\)

For \(i\in I\) , one denotes by \(\varepsilon_{i}\) the element of \(\mathbf{N}^{(I)}\) all of whose coordinates are zero, except for that of index \(i\) which is equal to 1. We have \(\alpha=\sum_{i\in I}\alpha_{i}\varepsilon_{i}\) for every element \(\alpha=(\alpha_{i})\) of \(\mathbf{N}^{(I)}\) .

